\documentclass[11pt,fontset=fandol]{ctexart}
\usepackage[a4paper,margin=23mm]{geometry}
\usepackage{amsmath,amssymb,mathtools,bm,mathrsfs,esint}
\usepackage{graphicx,longtable,booktabs,array,calc}
\usepackage[normalem]{ulem}
\usepackage{xcolor,hyperref}
\hypersetup{hidelinks}
\setcounter{secnumdepth}{-1}
\newcounter{none}
\providecommand{\tightlist}{\setlength{\itemsep}{0pt}\setlength{\parskip}{0pt}}
\providecommand{\passthrough}[1]{#1}
\usepackage{amssymb,mathrsfs,esint,bm,newunicodechar}
\allowdisplaybreaks
\setlength{\emergencystretch}{3em}
\providecommand{\llangle}{\langle\!\langle}
\providecommand{\rrangle}{\rangle\!\rangle}
\newunicodechar{∞}{\ensuremath{\infty}}
\newunicodechar{∈}{\ensuremath{\in}}
\newunicodechar{∑}{\ensuremath{\sum}}
\newunicodechar{∫}{\ensuremath{\int}}
\newunicodechar{≤}{\ensuremath{\leq}}
\newunicodechar{≥}{\ensuremath{\geq}}
\newunicodechar{≠}{\ensuremath{\neq}}
\newunicodechar{→}{\ensuremath{\to}}
\newunicodechar{⇒}{\ensuremath{\Rightarrow}}
\newunicodechar{α}{\ensuremath{\alpha}}
\newunicodechar{β}{\ensuremath{\beta}}
\newunicodechar{γ}{\ensuremath{\gamma}}
\newunicodechar{λ}{\ensuremath{\lambda}}
\newunicodechar{φ}{\ensuremath{\varphi}}
\newunicodechar{ε}{\ensuremath{\varepsilon}}
\newunicodechar{σ}{\ensuremath{\sigma}}
\newunicodechar{θ}{\ensuremath{\theta}}
\newunicodechar{δ}{\ensuremath{\delta}}
\newunicodechar{η}{\ensuremath{\eta}}
\newunicodechar{π}{\ensuremath{\pi}}
\newunicodechar{ξ}{\ensuremath{\xi}}
\newunicodechar{Σ}{\ensuremath{\Sigma}}
\newunicodechar{Δ}{\ensuremath{\Delta}}
\newunicodechar{∂}{\ensuremath{\partial}}
\newunicodechar{∀}{\ensuremath{\forall}}
\newunicodechar{∃}{\ensuremath{\exists}}
\newunicodechar{√}{\ensuremath{\surd}}
\newunicodechar{−}{\ensuremath{-}}
\newunicodechar{μ}{\ensuremath{\mu}}
\newunicodechar{ρ}{\ensuremath{\rho}}
\newunicodechar{τ}{\ensuremath{\tau}}
\newunicodechar{ω}{\ensuremath{\omega}}
\newunicodechar{Ω}{\ensuremath{\Omega}}
\newunicodechar{χ}{\ensuremath{\chi}}
\newunicodechar{ψ}{\ensuremath{\psi}}
\newunicodechar{Γ}{\ensuremath{\Gamma}}
\newunicodechar{Φ}{\ensuremath{\Phi}}
\newunicodechar{Λ}{\ensuremath{\Lambda}}
\newunicodechar{ℝ}{\ensuremath{\mathbb{R}}}
\newunicodechar{ℤ}{\ensuremath{\mathbb{Z}}}
\newunicodechar{ℕ}{\ensuremath{\mathbb{N}}}
\newunicodechar{△}{\ensuremath{\triangle}}
\newunicodechar{⊗}{\ensuremath{\otimes}}
\begin{document}
\section{2024年杭州电子科技大学601数学分析考研真题参考答案(宝哥版)}\label{ux5e74ux676dux5ddeux7535ux5b50ux79d1ux6280ux5927ux5b66601ux6570ux5b66ux5206ux6790ux8003ux7814ux771fux9898ux53c2ux8003ux7b54ux6848ux5b9dux54e5ux7248}

来源：\texttt{HDU/601数学分析/2024/2024年杭州电子科技大学601数学分析考研真题参考答案(宝哥版).pdf}。原
PDF 页序：1, 2, 3, 4。

\section{杭州电子科技大学2024年数学分析(回忆版)真题解析}\label{ux676dux5ddeux7535ux5b50ux79d1ux6280ux5927ux5b662024ux5e74ux6570ux5b66ux5206ux6790ux56deux5fc6ux7248ux771fux9898ux89e3ux6790}

\begin{quote}
依据原参考答案整理。极限展开及中断的和函数公式已由助手结合题面校正或补全；其余题号范围保持原文件内容。
\end{quote}

\subsection{一.填空题(8题，每题5分，共40分)}\label{ux4e00.ux586bux7a7aux98988ux9898ux6bcfux98985ux5206ux517140ux5206}

\begin{enumerate}
\def\labelenumi{\arabic{enumi}.}
\tightlist
\item
  \(\lim _ { n \to \infty } ( n ^ { 2 0 2 4 } + 2 ^ { n } ) ^ { \frac { 1 } { n } } =\)
\end{enumerate}

\subsection{【宝哥解答】}\label{ux5b9dux54e5ux89e3ux7b54}

\[
\lim _ { n \to \infty } ( n ^ { 2 0 2 4 } + 2 ^ { n } ) ^ { \frac { 1 } { n } } = \lim _ { n \to \infty } \left[ 2 ^ { n } ( { \frac { n ^ { 2 0 2 4 } } { 2 ^ { n } } } + 1 ) \right] ^ { \frac { 1 } { n } } = 2 \lim _ { n \to \infty } ( { \frac { n ^ { 2 0 2 4 } } { 2 ^ { n } } } + 1 ) ^ { \frac { 1 } { n } } = 2 ,
\]

这里，
\(\lim _ { n \to \infty } ( { \frac { n ^ { 2 0 2 4 } } { 2 ^ { n } } } + 1 ) = 1\)
\(\lim _ { n \to \infty } { \frac { 1 } { n } } = 0\) ，故
\(\lim _ { n \to \infty } ( { \frac { n ^ { 2 0 2 4 } } { 2 ^ { n } } } + 1 ) ^ { \frac { 1 } { n } } = 1 ^ { 0 } = 1\)

\begin{enumerate}
\def\labelenumi{\arabic{enumi}.}
\setcounter{enumi}{1}
\tightlist
\item
  \(\int \frac { x ^ { 3 } } { \sqrt { 1 - x ^ { 2 } } } d x =\) 盗版
\end{enumerate}

【宝哥解答】

\[
\begin{array} { l } { { \displaystyle { \int \frac { x ^ { 3 } } { \sqrt { 1 - x ^ { 2 } } } d x \stackrel { x = \sin t , t \in ( - \frac { \pi } { 2 } , \frac { \pi } { 2 } ) } { = } \int \frac { \sin ^ { 3 } t } { \cos t } \cos t d t = \int \sin ^ { 2 } t \sin t d t } } } \\ { { \displaystyle { = - \int ( 1 - \cos ^ { 2 } t ) d \cos t = \frac { 1 } { 3 } \cos ^ { 3 } t - \cos t + C \ ~ } } } \\ { { \displaystyle { = \frac { 1 } { 3 } ( 1 - x ^ { 2 } ) ^ { \frac { 3 } { 2 } } - \sqrt { 1 - x ^ { 2 } } + C } } } \end{array}
\]

\begin{enumerate}
\def\labelenumi{\arabic{enumi}.}
\setcounter{enumi}{2}
\tightlist
\item
  \(y = 2 ^ { x }\) 的麦克劳林级数的 \(x ^ { 2 0 2 4 }\) 的系数为
\end{enumerate}

\subsection{【宝哥解答】}\label{ux5b9dux54e5ux89e3ux7b54-1}

\(y = 2 ^ { x }\) 的麦克劳林级数的 \(x ^ { 2 0 2 4 }\) 的系数为
\({ \frac { y ^ { ( 2 0 2 4 ) } ( 0 ) } { 2 0 2 4 ! } } = { \frac { ( \ln 2 ) ^ { 2 0 2 4 } } { 2 0 2 4 ! } }\)
4.
\(\lim _ { n \to \infty } { \frac { ( n + 1 ) \sin x } { x n + 4 } }\)
的间断点为

【宝哥解答】

\[
\lim _ { n \to \infty } { \frac { ( n + 1 ) \sin x } { x n + 4 } } = { \left\{ \begin{array} { l l } { 0 , x = 0 } \\ { \displaystyle { \frac { \sin x } { x } } , x \neq 0 } \end{array} \right. } ,
\]

其间断点为 \(x = 0\)

5.求 \(x ^ { 2 } + y ^ { 2 } + z ^ { 2 } = 6\) 与
\(x ^ { 2 } + y ^ { 2 } = z\) 的交线的切线方程.

\subsection{【宝哥解答】}\label{ux5b9dux54e5ux89e3ux7b54-2}

\(\left\{ { \begin{array} { l } { x = { \sqrt { 2 } } \cos t } \\ { y = { \sqrt { 2 } } \sin t } \\ { z = 2 } \end{array} } \right.\)
联立方程
\(\left\{ { \begin{array} { l } { x ^ { 2 } + y ^ { 2 } + z ^ { 2 } = 6 } \\ { x ^ { 2 } + y ^ { 2 } = z } \end{array} } \right.\)
，可得 \(\left\{ x ^ { 2 } + y ^ { 2 } = 2 \right.\) 其参数方程为

任取交线上一点 \(( x _ { 0 } , y _ { 0 } , 2 )\) ，设参数为
\(t _ { 0 }\)
\(\left\{ \begin{array} { l l } { \displaystyle \frac { x - x _ { 0 } } { x ^ { \prime } ( t _ { 0 } ) } = \frac { y - y _ { 0 } } { y ^ { \prime } ( t _ { 0 } ) } } \\ { z = 2 } \end{array} \right.\)
即
\(\left\{ \begin{array} { l l } { \displaystyle { \frac { x - x _ { 0 } } { - y _ { 0 } } } = { \frac { y - y _ { 0 } } { x _ { 0 } } } } \\ { z = 2 } \end{array} \right.\)
，即
\(\left\{ { \begin{array} { l } { x _ { \mathrm { 0 } } x + y _ { \mathrm { 0 } } y = 2 } \\ { z = 2 } \end{array} } \right.\)

\subsection{6.没回忆到.}\label{ux6ca1ux56deux5fc6ux5230.}

\begin{enumerate}
\def\labelenumi{\arabic{enumi}.}
\setcounter{enumi}{6}
\tightlist
\item
  \(f ( x , y , z ) = x ^ { 2 } + x y + \cdots\) ，从
  \(P _ { 0 } ( - 1 , 0 , 2 )\) 到P(2,1,1)的方向导数.(没回忆全) 8.L为
  \(x ^ { 2 } + y ^ { 2 } = 1\) 上半圆周，则
  \(\int _ { L } { ( x ^ { 2 } + y ^ { 2 } ) d s } =\) 数学
\end{enumerate}

\subsection{【宝哥解答】}\label{ux5b9dux54e5ux89e3ux7b54-3}

L为 \(x ^ { 2 } + y ^ { 2 } = 1\) 上半圆周，故
\(\int _ { L } ( x ^ { 2 } + y ^ { 2 } ) d s = \int _ { L } d s = \pi\)

\subsection{二.计算题(共25分)}\label{ux4e8c.ux8ba1ux7b97ux9898ux517125ux5206}

1.(5分)
\(\lim _ { x \to 0 } { \frac { x - \sin { x } } { x ( e ^ { x } - 1 ) } }\)

【宝哥解答】

\[
\lim_{x\to0}\frac{x-\sin x}{x(e^x-1)}
=\lim_{x\to0}\frac{x^3/6+o(x^3)}{x^2+o(x^2)}=0.
\]

2.(10分)
\(\sum _ { n = 1 } ^ { \infty } ( - 1 ) ^ { n } { \frac { 2 n + 1 } { n } } x ^ { 2 n }\)
的收敛域与和函数.

【宝哥解答】

\(\lim _ { n \to \infty } \sqrt { \frac { 2 n + 1 } { 2 n + 3 } } = 1\)
数学 幂级数的收敛半径为 :在 x=±1处发散，故其收敛域为(-1,1).在(-1,1)内,

【助手整理补全：原识别稿在和函数处中断】在 \(|x|<1\)，

\[
\sum_{n=1}^{\infty}(-1)^n\frac{2n+1}{n}x^{2n}
=-\frac{2x^2}{1+x^2}-\ln(1+x^2).
\]

\subsection{三.证明题(每题10分，共30分)}\label{ux4e09.ux8bc1ux660eux9898ux6bcfux989810ux5206ux517130ux5206}

1.(10分) \(f ( x , y ) = \left| x - y \right| \varphi ( x , y )\)
\(\varphi ( x , y )\) 在(0,0)连续.

\begin{enumerate}
\def\labelenumi{(\arabic{enumi})}
\tightlist
\item
  \(f _ { x } ( 0 , 0 ) , f _ { y } ( 0 , 0 )\) 存在的充要条件；
\end{enumerate}

(2)f(x,y)在(0,0)的可微性.

\subsection{【宝哥解答】}\label{ux5b9dux54e5ux89e3ux7b54-4}

\begin{enumerate}
\def\labelenumi{(\arabic{enumi})}
\tightlist
\item
  \(f ( x , y ) = \left| x - y \right| \varphi ( x , y ) , \varphi ( x , y )\)
  在(0,0)连续.
\end{enumerate}

\[
\lim _ { x \to 0 ^ { + } } { \frac { f ( x , 0 ) - f ( 0 , 0 ) } { x } } = \lim _ { x \to 0 ^ { + } } { \frac { \left| x \right| \varphi ( x , 0 ) } { x } } = \lim _ { x \to 0 ^ { + } } { \frac { x \varphi ( x , 0 ) } { x } } = \varphi ( 0 , 0 ) \ ,
\]

\[
\lim _ { x \to 0 ^ { - } } { \frac { f ( x , 0 ) - f ( 0 , 0 ) } { x } } = \lim _ { x \to 0 ^ { - } } { \frac { \left| x \right| \varphi ( x , 0 ) } { x } } = \lim _ { x \to 0 ^ { - } } { \frac { - x \varphi ( x , 0 ) } { x } } = - \varphi ( 0 , 0 )
\]

\[
\lim _ { y \to 0 ^ { + } } \frac { f ( 0 , y ) - f ( 0 , 0 ) } { y } = \lim _ { y \to 0 ^ { + } } \frac { \left| y \right| \varphi ( 0 , y ) } { y } = \lim _ { y \to 0 ^ { + } } \frac { y \varphi ( 0 , y ) } { y } = \varphi ( 0 , 0 ) ,
\]

\[
\lim _ { y \to 0 ^ { - } } \frac { f ( 0 , y ) - f ( 0 , 0 ) } { y } = \lim _ { y \to 0 ^ { - } } \frac { \left| y \right| \varphi ( 0 , y ) } { y } = \lim _ { y \to 0 ^ { - } } \frac { - y \varphi ( 0 , y ) } { y } = - \varphi ( 0 , 0 ) .
\]

因此， \(f _ { x } ( 0 , 0 ) , f _ { y } ( 0 , 0 )\) 存在的充要条件是
\(\varphi ( 0 , 0 ) = - \varphi ( 0 , 0 )\) ，即
\(\varphi ( 0 , 0 ) = 0\)

\begin{enumerate}
\def\labelenumi{(\arabic{enumi})}
\setcounter{enumi}{1}
\tightlist
\item
  \(f ( x , y ) = \left| x - y \right| { \bigl [ } \varphi ( x , y ) - \varphi ( 0 , 0 ) { \bigr ] } + \left| x - y \right| \varphi ( 0 , 0 )\)
\end{enumerate}

对任意 \(( x , y ) \neq ( 0 , 0 )\)

\[
\frac { \| x - y \| [ \varphi ( x , y ) - \varphi ( 0 , 0 ) ] \| } { \sqrt { x ^ { 2 } + y ^ { 2 } } } \le \sqrt { 2 }\left|\varphi ( x , y ) - \varphi ( 0 , 0 )\right| \to 0 , ( x , y ) \to ( 0 , 0 ) \ ,
\]

故
\(\lim _ { ( x , y ) \to ( 0 , 0 ) } { \frac { \| x - y \| [ \varphi ( x , y ) - \varphi ( 0 , 0 ) ] \| } { \sqrt { x ^ { 2 } + y ^ { 2 } } } } = 0\)
，即
\(\lim _ { ( x , y ) \to ( 0 , 0 ) } { \frac { { \Big | } x - y { \Big | } { \big [ } \varphi ( x , y ) - \varphi ( 0 , 0 ) { \big ] } } { \sqrt { x ^ { 2 } + y ^ { 2 } } } } = 0\)

\[
\left| x - y \right| \left[ \varphi ( x , y ) - \varphi ( 0 , 0 ) \right] = 0 \cdot x + 0 \cdot y + o ( \sqrt { x ^ { 2 } + y ^ { 2 } } )
\]

\[
\big | x - y \big | \big [ \varphi ( x , y ) - \varphi ( 0 , 0 ) \big ]
\]

如果 \(\varphi ( 0 , 0 ) = 0\) ，则
\(\left| x - y \right| \varphi ( 0 , 0 )\) 在(0,0)处可微，且微分为0，
\(f ( x , y )\) 在原点可微。

如果 \(\varphi ( 0 , 0 ) \neq 0\) ，则
\(\left| x - y \right| \varphi ( 0 , 0 )\) 在(0,0)处两偏导数都不存在，
\(\left| x - y \right| \varphi ( 0 , 0 )\) 在原点不可微，故
\(f ( x , y )\) 在原点不可微。

综上所述， \(f ( x , y )\) 在原点的可微当且仅当
\(\varphi ( 0 , 0 ) = 0\)

【宝哥备注】

第(2)问见宝哥题解精粹第16讲2例14.
\end{document}
